\title{%
  How do you know it's secure?
  Passwords
}
\author{Daniel Bosk\thanks{%
    This material is available under the Creative Commons
    Attribution-NonCommercial-ShareAlike (CC-BY-NC-SA) 4.0 international
    license.
    The material was written with some aid from GitHub Copilot.
}}
\institute{%
  KTH EECS
}

\begin{frame}
  \maketitle
\end{frame}

\mode*

\begin{abstract}
  \input{abstract.tex}
\end{abstract}

\clearpage

\section{The papers}

\begin{frame}
  \begin{question}
    What part of security is not a science?
  \end{question}

  \mode<article>{%
    My take: one candidate is the rules we impose on users' passwords.
    When Komanduri et al.\ started, \textcquote{OfPasswordsAndPeople}{\textins*{e}ven
    the best current guidelines for designing password-composition
    policies \textelp{} are based on theoretical estimates \textelp{} or
    small-scale laboratory studies}.
    NIST computed what a rule \emph{allows}, not what users \emph{choose}
    under it.
    That makes the question a good opening: it is not that password
    policy cannot be a science, but that for a long time nobody had
    observed the thing the guidelines made claims about.
    The two papers are attempts to change that, and the seminar asks how
    well they succeed.
  }
\end{frame}

\begin{frame}
  Two attempts to put password policies on an empirical footing:
  \begin{itemize}
    \item \fullcite{OfPasswordsAndPeople}
    \item \fullcite{CanLongPasswordsBeSecureAndUsable}
  \end{itemize}
  \begin{onlyenv}<2>
    \begin{remark}
      Focus on the \emph{methods}, not the results!
      We return to \citetitle{WhyPhishingWorks} in the seminar on
      quantitative and qualitative methods.
    \end{remark}
  \end{onlyenv}
\end{frame}

\mode<article>{%
  In the vetcyb26h round, \citetitle{WhyPhishingWorks} was due with the
  quantitative and qualitative methods module.
  Some reflections included it, most did not, and two said that it was
  unclear which papers the reflection should cover.
  So this seminar focuses on the two password papers.
}

\section{From your reflections}

\begin{frame}
  \begin{block}{What most of you saw}
    \mode<presentation>{%
      \begin{itemize}
        \item Random assignment to policies: the policy is the intervention
        \item Entropy is a questionable operationalisation of strength;
          cracking is closer to an attacker
        \item MTurk, a hypothetical e-mail account, no stakes: weak external
          validity
        \item The cracker depends on its training data
        \item Several rules change at once (not full factorial)
        \item More drop-outs under the strictest policy
      \end{itemize}
    }
    \mode<article>{%
      Nearly all reflections identified the random assignment to policies
      as the intervention, questioned entropy as a measure of password
      strength and doubted the external validity of passwords created by
      MTurk workers for a hypothetical e-mail account.
      Several also noted that the cracker depends on its training data, that
      the conditions change several rules at once, and that drop-out was
      higher under the strictest policies.
      These points are worth only a brief confirmation; the seminar time is
      better spent on where the reflections disagree.
    }
  \end{block}

  \begin{onlyenv}<2>
    \begin{remark}
      The strongest reflections did not only list these problems.
      They used them to say \emph{what kind of answer} the papers can give.
    \end{remark}
  \end{onlyenv}
\end{frame}

\begin{frame}
  \begin{block}{Your open questions}
    \begin{itemize}
      \item Is it even an experiment? Where is the control group? Blinding?
      \item Is changing several rules at once bad science?
      \item Was there a hypothesis or not?
      \item Is internal validity an inductive or a deductive inference?
      \item Entropy and cracking disagree: which one do we trust?
      \item What does a difference of 10 bits mean in guesses?
      \item Would another cracker rank the policies differently?
      \item How do we balance control and realism?
    \end{itemize}
  \end{block}

  \mode<article>{%
    These questions are paraphrased from the reflections.
    The frames that follow take them one at a time.
  }
\end{frame}

\section{Experiments}

\begin{frame}
  \begin{question}
    Some of you wrote that these studies are not experiments: there is no
    control group, no hypothesis and no (double) blinding.
    Others called them controlled, even single-blind, experiments.
    Who is right?
  \end{question}

  \begin{onlyenv}<2>
    \begin{question}
      Look at the design itself.
      \begin{itemize}
        \item What is the intervention, and who decides it?
        \item What do we compare against, if not a control group?
        \item What \emph{can} be blinded when the treatment is a rule the
          participant must follow?
      \end{itemize}
    \end{question}
  \end{onlyenv}
\end{frame}

\begin{frame}
  \begin{example}[Komanduri et al.]
    From the methodology:
    \begin{itemize}
      \item Each participant was assigned
        \textcquote{OfPasswordsAndPeople}{randomly to one of five
        conditions}.
      \item No significant difference in age, gender or technical
        experience between the conditions.
      \item The consent form said the study was about
        \textcquote{OfPasswordsAndPeople}{behavior visiting secure
        websites}.
      \item basic8survey and basic8: same policy, different scenarios.
    \end{itemize}
  \end{example}
\end{frame}

\begin{frame}
  \begin{example}[Shay et al.]
    The paper says \textcquote{CanLongPasswordsBeSecureAndUsable}{We
    assigned participants to one of eight conditions}, but not how.
  \end{example}

  \begin{onlyenv}<2>
    \begin{question}
      Does it matter how?
    \end{question}
  \end{onlyenv}

  \mode<article>{%
    My take: yes, both are experiments.
    The intervention is the composition policy, and the researchers, not
    the participants, decide who gets which; that is what separates them
    from an observational study of leaked passwords, where each site's
    users chose their site.
    The missing pieces the reflections list (control group, hypothesis,
    blinding) are features some experiments have, not what makes one.
    Does it matter how Shay et al.\ assigned?
    Yes: if the assignment was not random, for example by time of sign-up,
    the conditions could differ in who signed up when, and the comparison
    would carry that difference.

    Two reflections looked for the randomisation in Shay et al.\ and did
    not find it; they were right, the paper does not state it.
    Neither paper has an untreated group: the conditions are compared
    pairwise with each other, with the familiar policy (basic8, and comp8
    in Shay et al.) as the one to compare against.
    Each participant was given one condition's scenario and policy, so
    they could not be blind to their own treatment; the consent form
    described the study as being about behaviour on secure websites.
    The pair basic8survey and basic8 varies only the scenario and keeps the
    policy fixed, which is how the authors examine the framing.
  }
\end{frame}

\begin{frame}
  \begin{question}
    comprehensive8 adds several requirements to basic8 at once: upper and
    lower case, a symbol, a digit and a dictionary check.
    Shay et al.\ write that they \textcquote{CanLongPasswordsBeSecureAndUsable}{did not employ a
    full-factorial study}.
    \begin{itemize}
      \item What can we conclude from such a condition, and what not?
      \item Is it bad science?
      \item How many conditions would a full-factorial design need?
    \end{itemize}
  \end{question}

  \mode<article>{%
    Shay et al.\ state the consequence themselves: \textcquote{CanLongPasswordsBeSecureAndUsable}{we
    are unable to evaluate the effect of changing each individual
    variable}, and they
    \textcquote{CanLongPasswordsBeSecureAndUsable}{might have missed interaction effects}.
    The conditions remain valid comparisons of \emph{policies} as
    deployed packages, which is what an administrator chooses between.

    \emph{What can we conclude?}
    That comprehensive8 as a whole gives passwords with more entropy than
    basic8; not which of its five requirements did it, nor whether two of
    them together do more than each alone.
    Note, though, that the designs are not without one-factor contrasts.
    In Komanduri et al., dictionary8 differs from basic8 only by the
    dictionary check, and basic8survey from basic8 only by the scenario.
    In Shay et al., basic12, basic16 and basic20 differ only in length:
    \textcquote{CanLongPasswordsBeSecureAndUsable}{This set of conditions
    varies only in length, so we can measure its impact on security and
    usability}.
    The same holds for 3class12 against 3class16 and 2word12 against
    2word16.
    The last pair even shows an interaction:
    \textcquote{CanLongPasswordsBeSecureAndUsable}{The 2word approach seems
    to be more effective at increasing password strength when combined
    with a length-16 requirement than with a length-12 requirement}.
    So the design isolates length, but not the class or word requirements.

    \emph{How many conditions would a full factorial need?}
    A full factorial crosses every level of every factor.
    My count of the factors in Shay et al.: minimum length (8, 12, 16, 20:
    four levels), character classes (none, three of four, all four: three
    levels), the two-word requirement (two levels) and the dictionary check
    (two levels).
    That is $4 \times 3 \times 2 \times 2 = 48$ conditions instead of 8.
    At the roughly thousand participants per condition that both papers
    used, that is about \num{48000} participants completing the study,
    against the \num{8143} who did; and Shay et al.\ needed \num{15108}
    starts to get those.

    \emph{Is it bad science?}
    My take: no, it is a choice of research question.
    One of you asked this, and the answer is that the study asks which of
    these policies, not which rule.
    It becomes bad science only if the conclusions claim more, say that
    length rather than the class requirement is what makes 3class16
    strong.
    The papers mostly avoid that; the useful exercise for the students is
    to find the sentences where a rule, not a policy, is credited.
  }
\end{frame}

\begin{frame}
  \begin{question}
    25\,\% dropped out under comprehensive8, compared to less than 19\,\%
    in the other conditions\autocite{OfPasswordsAndPeople}.
    Does random assignment still protect the comparison?
  \end{question}

  \mode<article>{%
    The drop-out happens after the assignment.
    Shay et al.\ name the risk: participants who drop out
    \textcquote{CanLongPasswordsBeSecureAndUsable}{might behave differently than those who do
    not, potentially biasing our results}.
    Komanduri et al.\ acknowledge \textcquote{OfPasswordsAndPeople}{a potential for selection bias}
    and argue that it is unlikely to change the relative entropy
    differences.
    In Shay et al., 83.0\,\% finished part one under comp8, against up to
    94.5\,\% in the other
    conditions\autocite{CanLongPasswordsBeSecureAndUsable}.
  }
\end{frame}

\section{Inferences}

\begin{frame}
  \begin{question}
    NIST predicted these entropies.
    What did Komanduri et al.\ find, and what kind of inference is that?
  \end{question}

  \begin{table}
    \begin{sidecaption}{%
      NIST's predicted entropy per policy and the entropy that Komanduri
      et al.\ estimated from 1000 passwords per condition
      \parencite[Table~1]{OfPasswordsAndPeople}.
    }[tab:nist-entropy]
      \begin{tabular}{lrr}
        \toprule
        Condition & NIST predicts & Measured\only<2->{ (underestimate)} \\
        \midrule
        basic8 & 18 bits & \only<2->{29.43 bits} \\
        dictionary8 & 24 bits & \only<2->{28.99 bits} \\
        basic16 & 30 bits & \only<2->{44.67 bits} \\
        comprehensive8 & 30 bits & \only<2->{34.30 bits} \\
        \bottomrule
      \end{tabular}
    \end{sidecaption}
  \end{table}
\end{frame}

\begin{frame}
  \begin{question}
    The authors say their estimates are
    \textcquote{OfPasswordsAndPeople}{significant underestimates of the true
    entropy}.
    \begin{itemize}
      \item Then what was refuted: NIST's prediction or the estimator?
      \item Does the answer differ for basic16 vs comprehensive8 and for
        dictionary8 vs basic8?
      \item And Shay et al.: which hypothesis did they test?
    \end{itemize}
  \end{question}

  \mode<article>{%
    The reflections split here.
    Some read Komanduri et al.\ as hypothetico-deductive: NIST predicts that
    basic16 and comprehensive8 give the same entropy, the measurement
    disagrees, so by modus tollens the prediction fails.
    Others saw no hypothesis at all.
    The paper supports the first reading: the dictionary8 and basic8
    conditions \textcquote{OfPasswordsAndPeople}{were designed to verify that hypothesis}, namely
    NIST's claim that a dictionary check adds entropy, and the authors
    state what they \textcquote{OfPasswordsAndPeople}{expected, a priori}.

    The estimator is an auxiliary hypothesis, as several of you noted.
    The authors argue that their estimates are underestimates at 1000
    samples, so they compare them only between conditions, with one
    exception: \textcquote{OfPasswordsAndPeople}{when we find our underestimates to
    be significantly higher than previously reported estimates}.
    So basic16 at 44.67 bits against NIST's 30 bits survives the auxiliary
    hypothesis: an underestimate that exceeds the prediction.
    The dictionary8 result is different: it is a comparison of two
    underestimates, and it rests on both being underestimated alike.

    Shay et al.\ state no hypothesis in advance; they grouped rules
    \textcquote{CanLongPasswordsBeSecureAndUsable}{in ways we hypothesized might balance
    usability and security}, which makes the study exploratory, as several of you
    wrote.
  }
\end{frame}

\begin{frame}
  \begin{question}
    One of you asked: random assignment makes the groups equal
    \enquote{only by chance}.
    Is the inference from \enquote{the groups differ} to \enquote{the
    policy caused it} then deductive or inductive?
  \end{question}

  \mode<article>{%
    My take: inductive.
    The conclusion can be false while the premises are true: a chance
    imbalance between the groups, or a difference that arises after the
    assignment (the drop-out above), would give the same observation
    without the policy causing it.
    Random assignment does not remove that possibility; it makes it
    possible to say how improbable it is.
    That is what Komanduri et al.'s permutation tests do for the entropy
    differences: they ask how often a difference at least as large would
    turn up if the condition labels were shuffled at random, that is, if
    the policy made no difference.
    A small number says that chance is an unlikely explanation, not an
    impossible one.

    One could also make it deductive by adding premises: if the groups
    differ only in the policy, and the measurements are the same in all
    groups, then a difference in the outcome comes from the policy.
    But then the uncertainty has only moved into those premises, which
    are auxiliary hypotheses again, as in the previous question.

    The student who asked suggested that internal validity is the
    inductive step about chance and external validity a further, larger
    inductive step to the population.
    I agree: both are inductive, and they generalise over different
    things, the first from this assignment to the causal claim for these
    participants, the second from these participants and this setting to
    others.
  }
\end{frame}

\section{Measurements}

\begin{frame}
  \begin{question}
    Entropy and cracking were measured on the same passwords.
    Do they agree?
  \end{question}

  \begin{table}
    \begin{sidecaption}{%
      Estimated entropy and passwords cracked by John the Ripper in
      heuristic mode, per condition
      \parencite[Tables~1 and~2]{OfPasswordsAndPeople}.
    }[tab:entropy-cracking]
      \begin{tabular}{lrr}
        \toprule
        Condition & Entropy & Cracked \\
        \midrule
        basic8 & 29.43 bits & \only<2->{188 of 972} \\
        dictionary8 & 28.99 bits & \only<2->{31 of 963} \\
        basic16 & 44.67 bits & \only<2->{9 of 971} \\
        comprehensive8 & 34.30 bits & \only<2->{0 of 997} \\
        \bottomrule
      \end{tabular}
    \end{sidecaption}
  \end{table}
\end{frame}

\begin{frame}
  \begin{question}
    Some of you saw a failed convergent-validity check: two
    operationalisations of strength that disagree.
    Others suggested that they measure different things.
    Which is it?
  \end{question}

  \begin{onlyenv}<2>
    \begin{remark}
      Entropy describes the whole distribution.
      A distribution with high entropy can still contain a small set of
      easily guessed passwords\autocite{OfPasswordsAndPeople}.
    \end{remark}
  \end{onlyenv}

  \mode<article>{%
    Komanduri et al.\ make this point themselves: \textcquote{OfPasswordsAndPeople}{it is
    possible even within a password distribution of high entropy to have a
    small set of passwords that are \enquote{easily} guessable}.
    Bonneau proves the general version: for any margin, there are
    distributions whose resistance to an attacker who needs to guess only
    a fraction of the passwords is lower than the Shannon entropy by more
    than that margin\autocite[Theorem~3.3.1]{GuessingHumanChosenSecrets2012}.
    So the disagreement need not mean that either instrument failed; they
    operationalise different properties, the predictability of the whole
    distribution and the weakness of its most guessable part.
    The comparison basic16 vs comprehensive8 has a further problem: John
    the Ripper \textcquote{OfPasswordsAndPeople}{is optimized for short passwords}, so the authors
    \textcquote{OfPasswordsAndPeople}{cannot say confidently how the two
    conditions compare for cracking protection}.
  }
\end{frame}

\begin{frame}
  \begin{question}
    basic16 has about 10 bits more entropy than comprehensive8.
    Is it then about $2^{10} \approx 1000$ times harder to guess?
  \end{question}

  \mode<article>{%
    One reflection asked this and noticed that it holds only if all
    passwords are equally likely.
    Bonneau shows that for an attacker who must guess values one at a
    time, \textcquote[§3.1.1]{GuessingHumanChosenSecrets2012}{Shannon
    entropy has no direct correlation to guessing difficulty}.
    The cracking results point the same way: John the Ripper cracked 9
    basic16 passwords and none under
    comprehensive8\autocite[Table~2]{OfPasswordsAndPeople}, though the
    authors warn that the tool favours short passwords.
  }
\end{frame}

\begin{frame}
  \begin{question}
    Shay et al.\ used a cracker adapted to long passwords.
    Which policy is stronger: basic16 or comp8?
  \end{question}

  \begin{onlyenv}<2>
    \begin{example}
      The share of passwords cracked depends on the guessing budget.
      \begin{description}
        \item[After $10^6$ guesses] the three basic conditions have
          significantly more passwords cracked than comp8.
        \item[After about $1.16\times 10^{12}$ guesses] comp8 is among the
          three weakest conditions.
      \end{description}
    \end{example}
  \end{onlyenv}

  \begin{onlyenv}<3>
    \begin{question}
      Back to the research question.
      \begin{itemize}
        \item So what is \enquote{how easily can we guess}?
        \item Would another cracker, or other training data, rank the
          policies differently?
      \end{itemize}
    \end{question}
  \end{onlyenv}

  \mode<article>{%
    The two orderings are from the security results of
    \textcite{CanLongPasswordsBeSecureAndUsable}.
    Strength is a curve over the attacker's guessing budget, not one
    number.
    The authors also note that they \textcquote{CanLongPasswordsBeSecureAndUsable}{had more data for
    some conditions than others}, so that a novel policy may look strong partly
    because the attacker lacks training data for it, an advantage that
    \textcquote{CanLongPasswordsBeSecureAndUsable}{might be temporary}.
    The cracker is an instrument whose readings depend on a theory of the
    attacker: the observation is theory-dependent.
  }
\end{frame}

\section{The research questions}

\begin{frame}
  \begin{itemize}
    \item How do different password policies affect the password
      distribution?
    \item How easily can we guess the passwords under different password
      policies?
  \end{itemize}

  \begin{question}
    Read the research questions closely.
    \begin{itemize}
      \item What is \enquote{the password distribution}? Can we observe it?
      \item Are the papers' answers absolute or relative?
    \end{itemize}
  \end{question}

  \mode<article>{%
    The reflections read \enquote{distribution} as the character
    distribution, as collisions between passwords, or as common patterns.
    None of the papers observes the distribution itself, only a sample of
    about a thousand passwords per condition, summarised by entropy, by
    composition statistics and by cracking curves.
    The strongest reflections concluded that the answers are
    \emph{comparative} (one policy against another), \emph{conditional}
    (on this attacker and this training data) and tied to \emph{one
    population} (MTurk workers creating a password for a study).
    Komanduri et al.\ make the same argument for their ecological validity:
    since the hypothetical scenario was the same in the four conditions,
    they are \textcquote{OfPasswordsAndPeople}{confident that our comparisons among
    conditions are valid}.
  }
\end{frame}

\begin{frame}
  \begin{activity}[Design]
    Design a study that answers \emph{exactly} the two research questions.
    \begin{itemize}
      \item What would you need to observe?
      \item Real passwords from a deployed system, leaked password sets, or
        an experiment: what does each give you, and what does it cost?
    \end{itemize}
  \end{activity}

  \mode<article>{%
    Komanduri et al.\ describe the ideal and why they did not reach it:
    data from real users of a deployed system, which is hard because of
    the sensitivity of password data and the difficulty of partitioning
    real users into experimental conditions.
    One reflection proposed comparing study passwords with leaked
    passwords created under the same policy.

    My take, and it is also the answer to \enquote{how do we balance
    control and realism?}: no single source answers both research
    questions, so combine them.
    \begin{description}
      \item[The experiment] gives the causal comparison between policies,
        because the researchers assign them, but its passwords protect
        nothing.
      \item[Leaked sets] give real passwords protecting real accounts, but
        each site's policy comes with its own users, so a difference
        between sites is not a difference between policies.
      \item[A check between them] tells us how far the experiment can be
        trusted: Shay et al.\ cite two studies that compared study
        passwords with real ones; \textcquote{CanLongPasswordsBeSecureAndUsable}{Both
        studies concluded that passwords created during studies can
        resemble real, high-value passwords, yet are not a perfect
        proxy}.
    \end{description}
    For the second research question the design must also fix the
    attacker: online or offline, how many guesses, what training data.
    Without that, \enquote{how easily} has no answer, as Shay et al.'s
    reversed ranking showed.
  }
\end{frame}

\section{Further}

\begin{frame}
  \begin{question}
    We had \citetitle{OfPasswordsAndPeople}, why do we need
    \citetitle{CanLongPasswordsBeSecureAndUsable}?
  \end{question}

  \mode<article>{%
    My take: because the first paper's instrument could not answer the
    question it raised.
    Komanduri et al.\ found basic16 to have the most entropy but
    \textcquote{OfPasswordsAndPeople}{cannot say confidently how the two
    conditions compare for cracking protection}, since John the Ripper
    favours short passwords.
    Shay et al.\ name what the entropy figure hid: under the length-16
    policy \textcquote{CanLongPasswordsBeSecureAndUsable}{there were also
    very simple passwords, like passwordpassword, that could be guessed
    easily}.
    So the second paper brings a better instrument (a cracker adapted to
    long passwords, read over guessing budgets), new conditions aimed at
    removing that weak tail, and a second sample in which the basic16 and
    comp8 comparison is made again.
    The last point is replication, which several of you mentioned: one
    study's ranking is not yet a result.
  }

  \begin{question}
    What factors affect how users choose passwords?
    \begin{itemize}
      \item Is it just the password policy?
      \item Can we know from these studies?
      %\item We'd need more studies focusing on other factors, then a
        %meta-analysis to see what factors are most important and how they
        %interact.
    \end{itemize}
  \end{question}

  \mode<article>{%
    My take: not just the policy, and the papers can tell us about some
    of the other factors with very different confidence.
    \begin{description}
      \item[The scenario] was randomised in Komanduri et al., so here we
        can make a causal claim: the four e-mail conditions gave
        significantly more entropy than basic8survey, which the authors
        read as \textcquote{OfPasswordsAndPeople}{consistent with users
        picking better passwords to protect a hypothetical email account
        than to protect a real survey account}.
      \item[Writing the password down] went with an extra 1.9 bits of
        entropy\autocite{OfPasswordsAndPeople}.
        But participants chose whether to write it down, so this is a
        correlation: it may be that people who make long passwords also
        write them down, not that writing down causes long passwords.
      \item[Reusing an existing password] was reported more often under
        the less demanding policies\autocite{OfPasswordsAndPeople}.
        Here the policy is randomised, so the policy affects reuse; but
        what reuse then does to the password is again only observed.
    \end{description}
    So the general answer: we can know the effect of what was assigned
    (policy, scenario), and only associations for what participants
    chose themselves.
    That is the same distinction as between the experiment and the leaked
    sets in the design activity.
  }
\end{frame}

\begin{frame}
  \begin{question}
    Did they make use of any models?
  \end{question}

  \mode<article>{%
    My take: several, and each one is something the conclusions rest on.
    \begin{description}
      \item[NIST's estimates] are a model of what a policy yields,
        computed from what the rules allow; the experiment tests it.
      \item[The entropy estimate] is a model of the password
        distribution: Komanduri et al.\ estimate it by
        \textcquote{OfPasswordsAndPeople}{summing the entropy values
        derived from each element of a password}, which assumes that the
        elements contribute separately.
      \item[The cracker] is a model of the attacker; Shay et al.'s
        algorithm \textcquote{CanLongPasswordsBeSecureAndUsable}{uses a
        corpus of training data to generate guesses in order of
        likelihood, up to some cutoff}, so it is only as good as that
        corpus.
      \item[The participants] stand in for real users, and the e-mail
        scenario for a real account: a model in the sense that one thing
        is studied in place of another.
    \end{description}
  }
\end{frame}

\begin{frame}
  \begin{example}
    Entropy as a model for password strength (guessability).
    We saw above that it does not capture the weakest passwords.
  \end{example}

  \begin{remark}
    This is the topic of the next section of the course.
  \end{remark}
\end{frame}

